\section{Principios de ingenieria y arquitectura}

\subsection{Descripcion}

Este documento fija los principios que deben reflejarse en la documentacion y, posteriormente, en la implementacion del sistema. No es una lista decorativa; actua como criterio de calidad para escribir reglas, contratos y limites.

\subsection{Alcance}

Aplica a backend, frontend, datos, integraciones y operacion.

\subsection{Definiciones}

\begin{itemize}
  \item \textbf{KISS}: la solucion documental y tecnica debe ser lo mas simple posible sin sacrificar claridad.
  \item \textbf{Single Responsibility}: cada documento y cada modulo deben tener una responsabilidad dominante.
  \item \textbf{Explicit Contracts}: toda frontera debe declarar que recibe, que entrega y bajo que reglas opera.
\end{itemize}

\subsection{Reglas}

Principios de ingenieria aplicados al repositorio:

\begin{itemize}
  \item `KISS`: los documentos deben ser directos y libres de teoria innecesaria.
  \item `SOLID` y `Single Responsibility`: un documento responde una sola pregunta principal.
  \item `DRY`: una regla normativa se define una vez y luego se referencia.
  \item `YAGNI`: no se documentan capacidades fuera del alcance confirmado.
  \item `Separation of Concerns`: producto, dominio, arquitectura, datos, seguridad y operacion viven en documentos separados.
  \item `Law of Demeter`: se evita acoplar documentos a detalles de otras capas cuando no son necesarios.
  \item `Fail Fast`: se documentan estados invalidos, entradas invalidas y errores esperados.
  \item `Least Privilege`: los roles se definen con permisos minimos suficientes.
  \item `Idempotency`: recalculo, sincronizacion y acciones repetibles deben quedar documentadas como operaciones seguras frente a reintentos.
  \item `Convention over Configuration`: el repositorio adopta convenciones fijas para estructura, nombres y compilacion.
  \item `Explicit Contracts`: toda interfaz relevante debe quedar escrita de forma verificable.
\end{itemize}

Principios arquitectonicos confirmados:

\begin{itemize}
  \item Frontend organizado por features.
  \item Separacion entre UI, logica y acceso a datos.
  \item `shared` restringido a elementos realmente reutilizables.
  \item Backend como modular monolith.
  \item Separacion por dominio: auth, users, payments, matches, picks, results, scoring, leaderboard, admin, audit y notifications.
  \item Validacion en la frontera de entrada.
  \item Logica de negocio fuera de controllers.
\end{itemize}

Lineamientos de produccion:

\begin{itemize}
  \item `Secure by Design`.
  \item Buenas practicas `OWASP`.
  \item `Twelve-Factor App`.
  \item Configuracion fuera del codigo.
  \item Logs como stream de eventos.
  \item Validacion estricta de inputs.
  \item Manejo claro de errores.
  \item Auditoria de acciones criticas.
\end{itemize}

Baseline aprobado:

\begin{itemize}
  \item Frontend: Next.js, TypeScript, ESLint, Zod, TanStack Query, Tailwind CSS, Playwright.
  \item Backend: Node.js, Express, TypeScript, Zod, Prisma, Helmet, Pino, OpenTelemetry, Vitest.
  \item Base de datos: PostgreSQL.
\end{itemize}

\subsection{Casos borde}

\begin{itemize}
  \item Si una practica recomendada entra en conflicto con una regla del dominio, la regla del dominio no se altera sin decision explicita.
  \item Si una herramienta baseline cambia en el futuro, debe registrarse como decision del repositorio.
\end{itemize}

\subsection{Invariantes}

\begin{itemize}
  \item La documentacion debe justificar limites y responsabilidades; no solo listarlos.
  \item Ningun contrato relevante debe quedar implicito.
  \item Las decisiones tecnicas deben preservar claridad operativa y trazabilidad.
\end{itemize}

\subsection{Preguntas abiertas}

\begin{itemize}
  \item Ninguna por ahora.
\end{itemize}

